\clearpage
\section{Exact enlarged-carrier recovery and comparison with the original contract}\label{app:lowrank_contract}

Evidence: a bounded contract audit and elementary conditional recovery
lemmas. This note does not prove that the proposed compact enlargement
is a group, that its induced representation is Heisenberg, or that the
source markings attach to it. Those are the arithmetic inputs to be
verified by the separate rank-two and rank-four constructions.

\textbf{Verdict:} an enlarged compact carrier is not, merely by being
Heisenberg or an equivalent type, a new instance of the original H2
quotient. It can serve as a common assembly carrier for the already
supplied H2 targets if its actual occurrence, marked attachment, and
same-candidate recovery are proved and retained. The recovery mechanism
differs from the finite sector idempotent used when the packet lies
inside the original inducing carrier.

\hypertarget{lowrank_contract-what-original-h2-literally-requires}{%
\subsection{1. What original H2 literally
requires}\label{lowrank_contract-what-original-h2-literally-requires}}

Original Section~10 fixes an actual maximal simple character theta, its
actual groups H\_theta¹ subset J\_theta¹ subset J\_theta, and

\begin{Verbatim}[fontsize=\small,breaklines=true,breakanywhere=true]
V_theta=J_theta¹/H_theta¹.
\end{Verbatim}

Its Heisenberg representation is the particular eta\_theta containing
theta. Definition~11.1 supplies an equivariant symplectic isometry from a
source envelope to that type-side quotient, together with the central
character, and at p=2 the compatible quadratic refinements and
central-extension lift. Standard Fact 10.2 separately supplies the
canonical beta extension of this simple-character datum.

Consequently the following implications are invalid without more proof:

\begin{Verbatim}[fontsize=\small,breaklines=true,breakanywhere=true]
Xi is a finite Heisenberg representation
   therefore Xi=eta_theta' for a maximal simple character theta';

(P,Xi) is equivalent to a known arithmetic type
   therefore it has that type's J¹/H¹ or canonical beta extension.
\end{Verbatim}

An equivalence of generated representation categories does not identify
the particular finite group quotient in Definition~11.1. If a fixed
conductor and torsion profile forces every maximal simple-character
carrier in the candidate to have dimension one or two, constructing a
larger Xi in that candidate does not refute the dimension theorem. Xi is
another arithmetic carrier.

There is nevertheless an important distinction between \textbf{replacing
an H2 target} and \textbf{assembling already supplied H2 targets}. The
accepted sector construction itself uses a finite sector that need not
be a whole maximal simple-character carrier. Its old actual break
targets remain unchanged and are attached to that sector. An enlarged
P-carrier can play the analogous assembly role, provided the complete
dictionary and recovery below are proved. Calling it merely an arbitrary
vector space would understate that result. Calling it a new canonical maximal simple type would overstate it.

\hypertarget{lowrank_contract-precise-conditional-setup}{%
\subsection{2. Precise conditional
setup}\label{lowrank_contract-precise-conditional-setup}}

Let G=GL\_n(F), Z=F×I\_n, and let the already constructed candidate be

\begin{Verbatim}[fontsize=\small,breaklines=true,breakanywhere=true]
pi=c-Ind_N^G Lambda,
\end{Verbatim}

with its fixed actual extended type and fixed central character omega.
Let J be a compact open normal subgroup of N such that

\begin{Verbatim}[fontsize=\small,breaklines=true,breakanywhere=true]
eta=Lambda|J is irreducible,       N/J is infinite cyclic.
\end{Verbatim}

In the proposed F2-residue, full-degree totally ramified setting, J is
the actual compact type group J¹=J, not a new group called J merely to
fit this notation. Let P be an actual compact open subgroup of G
satisfying

\begin{Verbatim}[fontsize=\small,breaklines=true,breakanywhere=true]
J subset P,          N intersect P=J.
\end{Verbatim}

Define the genuine finite-dimensional compact induction

\begin{Verbatim}[fontsize=\small,breaklines=true,breakanywhere=true]
Xi=Ind_J^P eta.                                       (1)
\end{Verbatim}

Neither irreducibility nor Heisenberg structure of Xi is assumed to
follow from its dimension. They must be proved separately when needed.

If the known eta intertwining satisfies I\_G(eta)=N, then Xi is
irreducible: the finite Mackey endomorphism formula gives

\begin{Verbatim}[fontsize=\small,breaklines=true,breakanywhere=true]
End_P(Ind_J^P eta)=End_J(eta)=C,
\end{Verbatim}

because no nontrivial double coset in P lies in I\_G(eta).
Finite-dimensional smooth representations of compact P factor through a
finite group and are semisimple, so scalar endomorphisms imply
irreducibility. This criterion requires its actual group and intertwining hypotheses. Dimension alone is not a substitute.

\hypertarget{lowrank_contract-xi-occurs-in-the-same-pi-and-can-be-an-equivalent-compact-type}{%
\subsection{3. Xi occurs in the same pi, and can be an equivalent
compact
type}\label{lowrank_contract-xi-occurs-in-the-same-pi-and-can-be-an-equivalent-compact-type}}

The subspace of the actual compact-induction model of pi consisting of
functions supported on NP is stable under P. Mackey restriction
identifies that subspace with

\begin{Verbatim}[fontsize=\small,breaklines=true,breakanywhere=true]
Ind_(N intersect P)^P(Lambda|_(N intersect P))=Xi.     (2)
\end{Verbatim}

Thus Xi is an actual compact P-subrepresentation of the same pi, with no
change of arithmetic candidate or noncompact scalar.

Induction in stages also gives the exact G-isomorphism

\begin{Verbatim}[fontsize=\small,breaklines=true,breakanywhere=true]
c-Ind_P^G Xi isomorphic to c-Ind_J^G eta.             (3)
\end{Verbatim}

If (J,eta) is already a type for a stated Bernstein block and Xi is
irreducible, then (P,Xi) detects the same irreducible representations:

\begin{Verbatim}[fontsize=\small,breaklines=true,breakanywhere=true]
Hom_P(Xi,V) isomorphic to Hom_J(eta,V).
\end{Verbatim}

This is the precise statement about equivalent types that follows from (1). It
does not assert that Xi is the eta of another maximal simple character,
or attach a canonical beta extension to Xi. The known canonical
extension in Lambda remains the one originally supplied.

\hypertarget{lowrank_contract-exact-recovery-before-fixing-the-center-a-quotient}{%
\subsection{4. Exact recovery before fixing the center: a
quotient}\label{lowrank_contract-exact-recovery-before-fixing-the-center-a-quotient}}

Choose u in N mapping to a generator of N/J, and retain the actual
matrix Lambda(u). Use the model

\begin{Verbatim}[fontsize=\small,breaklines=true,breakanywhere=true]
X=c-Ind_J^G eta,
f(jg)=eta(j)f(g),
\end{Verbatim}

with G acting by right translation. Define the invertible G-equivariant
shift operator

\begin{Verbatim}[fontsize=\small,breaklines=true,breakanywhere=true]
(Tf)(g)=Lambda(u)^(-1)f(ug).                          (4)
\end{Verbatim}

Normality of J and the actual intertwining relation for Lambda(u) ensure
that Tf has the same J-covariance. Its normalization explicitly retains the fixed noncompact extension. Changing Lambda(u) changes the relation
below.

Define

\begin{Verbatim}[fontsize=\small,breaklines=true,breakanywhere=true]
(Qf)(g)=sum_(k in Z) Lambda(u)^(-k)f(u^k g).         (5)
\end{Verbatim}

The sum is finite on each argument and has the required compact support
modulo N. It is N-covariant and therefore lies in pi. It is surjective
and satisfies Q(Tf)=Q(f). Decomposing support into N/J-orbits reduces
its kernel to the elementary augmentation kernel on finitely supported
sequences on Z. That kernel consists exactly of consecutive differences.
Hence

\begin{Verbatim}[fontsize=\small,breaklines=true,breakanywhere=true]
X/(T-1)X isomorphic to pi.                           (6)
\end{Verbatim}

Through (3), this is a recovery of pi from the enlarged compact carrier
plus the original noncompact relation. It does not select the correct
relation from the finite Heisenberg packet alone.

In general this quotient cannot be replaced by a G-equivariant split
inclusion before central normalization. A nonzero vector of X has finite
support modulo compact J. The valuation of its determinant support is a
finite subset of Z. A scalar uniformizer z=tI\_n shifts these valuations by n. Consequently X has no nonzero eigenvector for z. Every vector of
pi is such an eigenvector, with eigenvalue omega(t). Therefore no
nonzero G-map pi -\textgreater{} X exists. In particular (6) does not
split.

This is why the Block159 sector summand projector cannot simply be
quoted: there the finite sector was inside Lambda and its stabilizer
contained the noncompact center. Here Xi is larger than Lambda and lies
in pi\textbar P, with P compact. These are different induction problems.

\hypertarget{lowrank_contract-after-retaining-the-center-a-finite-hecke-idempotent}{%
\subsection{5. After retaining the center: a finite Hecke
idempotent}\label{lowrank_contract-after-retaining-the-center-a-finite-hecke-idempotent}}

There is also an exact summand version once the already fixed central
character is imposed. Put J\_Z=ZJ and P\_Z=ZP. Assume, as in the
proposed totally ramified full-degree setting, that

\begin{Verbatim}[fontsize=\small,breaklines=true,breakanywhere=true]
N/J_Z is finite cyclic of order ell
\end{Verbatim}

(here ell=n). Let Lambda\_0=Lambda\textbar J\_Z. Because Xi occurs in
the same pi, its action on Z intersect P agrees with omega. It therefore
extends to P\_Z by letting Z act through omega. Call this Xi\_omega.
Induction in stages gives

\begin{Verbatim}[fontsize=\small,breaklines=true,breakanywhere=true]
Y=c-Ind_(P_Z)^G Xi_omega
  isomorphic to c-Ind_(J_Z)^G Lambda_0.              (7)
\end{Verbatim}

Using the same normalized shift (4), now on Y, one has T\^{}ell=1.
Indeed u\^{}ell belongs to J\_Z and the J\_Z covariance cancels
Lambda(u)\^{}ell. The finite operator

\begin{Verbatim}[fontsize=\small,breaklines=true,breakanywhere=true]
e_Lambda=(1/ell) sum_(j=0)^(ell-1) T^j               (8)
\end{Verbatim}

is a G-equivariant idempotent whose image is exactly the fixed pi. To
see this directly, N-covariant functions embed in J\_Z-covariant
functions because N/J\_Z is finite. On that embedded pi the finite
version of (5) equals ell times the identity. Its average is precisely
(8), and its image consists exactly of the N-covariant functions.

Equivalently, Ind\_(J\_Z)\^{}N Lambda\_0 is Lambda tensored with the
regular representation of N/J\_Z. Formula (8) selects the Lambda summand.
It uses the already fixed Lambda(u), including its root choice, so it
preserves the same candidate. Other root choices select the other
extensions in this finite induction. They are not interchangeable input
data.

This is a finite Hecke/induction projector after normalization of the central character. It is not the original finite projector in the image algebra onto a sector of eta. A final artifact must retain (7), (8), and the
normalization Lambda(u), rather than describing recovery merely as ``the
carrier is an equivalent type.''

\hypertarget{lowrank_contract-the-exact-remaining-admissibility-checks}{%
\subsection{6. The exact remaining admissibility
checks}\label{lowrank_contract-the-exact-remaining-admissibility-checks}}

For either proposed low-rank enlargement, the new proof must supply all
of the following concrete inputs:

\begin{enumerate}
\def\labelenumi{\arabic{enumi}.}
\tightlist
\item
  \textbf{Actual compact group:} prove that the proposed P is a
  subgroup, is compact open, contains the actual J, and satisfies P
  intersect N=J. Compute {[}P:J{]} and the resulting Xi dimension.
\item
  \textbf{Actual carrier:} verify eta=Lambda\textbar J, its actual
  intertwining if the irreducibility criterion is used, and the Mackey
  realization (2).
\item
  \textbf{Actual finite Heisenberg structure:} exhibit the scalar
  preimage in P and the resulting elementary abelian projective
  quotient. Compute its represented commutator and prove nondegeneracy.
  An irreducible carrier of dimension two or four need not automatically
  have this structure.
\item
  \textbf{Original attachments:} attach the unchanged supplied actual H2
  break targets, retaining their full source character values and dyadic
  finite extension/lift/square markings. The precise distinction between transported and native markings remains the one proved in
  \path{DYADIC_MARKINGS167.md}.
\item
  \textbf{Recovery data:} retain the actual Lambda, central character
  and noncompact generator normalization, and prove the isomorphisms and
  quotient/idempotent above for those exact groups.
\end{enumerate}

When these are supplied, the result is stronger than occurrence of an
unrelated finite vector space: it is a marked arithmetic carrier in the
same candidate, equivalent through explicit induction and recovering
that candidate exactly. For C=\{1\} and one vertex, it can repair the
failed low-rank \textbf{assembly-carrier} realization under the retained
scalar-pushout contract. The original H2 targets remain their original simple-character quotients. They are not renamed Xi.

If the intended target instead requires the source packet itself to
equal J\_theta'¹/H\_theta'¹ for a maximal simple character in this same
candidate, the enlargement does not solve it. The retained low-rank
dimension obstruction still applies. A separate theorem identifying Xi
with such an eta would be needed, and under the fixed obstructed profile
it cannot be inferred from type equivalence.

No H5 condition is added or discharged. If H5 is subsequently supplied,
it must concern this fixed recovered candidate. Nontrivial stabilizer
actions, canonical native quadratic comparison, and arbitrary transport
graphs remain outside this bounded argument.

\hypertarget{lowrank_contract-evidence-and-self-review}{%
\subsection{Evidence and self-review}\label{lowrank_contract-evidence-and-self-review}}

Read directly: original Section~10 and Definition~11.1 in the raw
manuscript; the Block159 \path{NONINVARIANT_SECTOR_RECONSTRUCTION159.md};
current \path{PROOF167.md} and \path{TARGET167.md}; the retained dyadic marking
theorem. The Mackey, induction, shift-quotient, and central-normalized
projector arguments above are proved conditionally for the displayed
groups, rather than assumed for the proposed congruence formula. No
actual low-rank P computation or new Lean execution is claimed in this
contract note.
